% !TEX program = lualatex
% STOC 2027 author video (20--30 min, for PC members and subreviewers)
%   Story: main result -> why it is hard -> how we solve it (details are in the paper)
%   Rules: https://acm-stoc.org/stoc2027/stoc2027-video-instructions.pdf
%   Style: follows ../../PCORTC_derivative/lics2026_slide/slide.tex
\documentclass[unicode,11pt,aspectratio=169]{beamer}

\usetheme[numbering=fraction,sectionpage=progressbar,subsectionpage=progressbar]{moloch}
\molochset{block=transparent}
\setbeamercovered{invisible}
\setbeamertemplate{navigation symbols}{}
\setbeamertemplate{itemize item}{\textbullet}
\setbeamertemplate{itemize subitem}{--}
\setbeamertemplate{enumerate item}{\insertenumlabel.}
\setbeamertemplate{enumerate subitem}{\insertsubenumlabel.}
\setbeamertemplate{bibliography item}{}

\usefonttheme[onlymath]{serif}
\usepackage{fontspec}
\usepackage{amsmath,amssymb,mathtools,stmaryrd,amsthm}
\usepackage{color}
\usepackage{booktabs,array,colortbl}
\usepackage{graphicx}
\usepackage{tikz}
\usetikzlibrary{tikzmark}
\usetikzlibrary{arrows.meta,positioning,fit,calc,backgrounds,decorations.pathreplacing,shapes.callouts,shapes.geometric}

\usepackage[backend=biber, uniquelist=false, maxcitenames = 2, maxbibnames=8, style=authoryear-comp,sorting = none, firstinits=true, sortcites=false]{biblatex}
\DeclareNameAlias{sortname}{given-family}
\DeclareNameAlias{default}{given-family}
\addbibresource{references.bib}

% Plain, monochrome style: black text, gray structure, a single dark accent.
\definecolor{chiba}{RGB}{96,39,40}
\colorlet{chibasub}{chiba}
\colorlet{logicblue}{black}
\colorlet{logicgreen}{black}
\colorlet{logicorange}{black!55}

\setbeamercolor{normal text}{fg=black,bg=white}
\setbeamercolor{alerted text}{fg=chiba}
\setbeamercolor{title separator}{fg=black!60,bg=black!60}
\setbeamercolor{progress bar}{fg=black!60,bg=black!15}
\setbeamercolor{frametitle}{bg=white,fg=black}
\setbeamercolor{block title}{fg=black,bg=}
\setbeamercolor{block body}{fg=black,bg=}
\setbeamercolor{block title alerted}{fg=black,bg=}
\setbeamercolor{block body alerted}{fg=black,bg=}
\setbeamercolor{block title example}{fg=black,bg=}
\setbeamercolor{block body example}{fg=black,bg=}
\setbeamerfont{block title}{series=\bfseries}
\setbeamertemplate{frametitle}{%
  \vspace*{.6em}\hspace*{-.2em}{\usebeamerfont{frametitle}\bfseries\insertframetitle}\par
  \vspace*{-.3em}{\color{black!40}\rule{\linewidth}{.4pt}}\par}

\makeatletter
\setlength{\moloch@progressonsectionpage@linewidth}{2pt}
\makeatother

\newtheorem{proposition}[theorem]{Proposition}

\AtBeginSection[]{}

\AtBeginDocument{
  \abovedisplayskip=.5\abovedisplayskip
  \abovedisplayshortskip=.5\abovedisplayshortskip
  \belowdisplayskip=.5\belowdisplayskip
  \belowdisplayshortskip=.5\belowdisplayshortskip
}

\renewcommand{\emph}[1]{\textbf{\textup{#1}}}
\newcommand{\EMPH}[1]{\textbf{#1}}
\newcommand{\KEY}[1]{\textbf{#1}}
\newcommand{\NOTE}[1]{{\footnotesize\textcolor{black!55}{#1}}}
\newcommand{\fcite}[2][]{{\textcolor{gray!50!black!80}{\footnotesize(\cite[#1]{#2})}}}
\newcommand{\HP}[1]{$\Rightarrow$\quad
  \begin{minipage}[t]{\dimexpr\linewidth-2.2em\relax}#1\end{minipage}}

% ==== notation (same as the paper ../main.tex) ====
\newcommand{\defeq}{\coloneqq}
\newcommand{\set}[1]{\{#1\}}
\newcommand{\CoR}{\ensuremath{\mathrm{CoR}}}
\newcommand{\FOthree}{\ensuremath{\mathrm{FO}^{3}}}
\newcommand{\FOfour}{\ensuremath{\mathrm{FO}^{4}}}
\newcommand{\zero}{\mathbf{0}}
\newcommand{\one}{\mathbf{1}}
\newcommand{\id}{\mathbf{1}'}
\newcommand{\conv}{\smile}
\newcommand{\cost}[1]{\#_{;}(#1)}
\newcommand{\gcost}[1]{g_{;}(#1)}
\newcommand{\wid}[1]{\mathsf{sn}(#1)}
\newcommand{\Col}{\Lambda}
\newcommand{\colr}{\chi}
\newcommand{\cell}[1]{C_{#1}}
\newcommand{\ctype}[1]{\tau_{#1}}
\newcommand{\tpl}[1]{\tau_{#1,\ast}}
\newcommand{\tpr}[1]{\tau_{\ast,#1}}
\newcommand{\tpeq}{\tau_{=}}
\newcommand{\tp}{\operatorname{tp}}
\newcommand{\res}{\mathord{\restriction}}
\newcommand{\M}{\mathfrak M}
\newcommand{\eqfin}{\mathrel{\equiv_{\mathrm{fin}}}}
\newcommand{\cbinom}{\tbinom{k}{\lfloor k/2 \rfloor}}

% colours of the three cells in the figure of M
\colorlet{colA}{logicblue}
\colorlet{colB}{logicorange}
\colorlet{colC}{logicgreen}

\tikzset{
  cellvert/.style={circle,fill=black!75,inner sep=0pt,minimum size=4pt},
  anchor vert/.style={circle,draw=black!60,fill=white,thick,inner sep=0pt,minimum size=8pt},
  hubvert/.style={circle,fill=black,inner sep=0pt,minimum size=10pt},
  cellshape/.style={ellipse,draw=black!45,fill=black!3,minimum width=2.6cm,minimum height=1.7cm},
  stepbox/.style={draw=#1,thick,fill=#1!6,rounded corners=2pt,align=left,
    text width=4.15cm,minimum height=2.6cm,inner sep=6pt,anchor=north},
  ev/.style={circle,fill=#1,draw=white,thick,minimum size=.32cm,inner sep=0pt},
  lbl/.style={align=left,font=\small,anchor=west,text width=11.5cm},
  logicrow/.style={draw=black,thin,rounded corners=2pt,minimum height=.75cm,minimum width=3.2cm,align=center},
}

\hypersetup{colorlinks=false}

\title{An Exponential Succinctness Gap between\\Three-Variable Logic and the Calculus of Relations}
\author{Yuya Uezato \and Yoshiki Nakamura}
\institute{CyberAgent, Inc. \qquad Chiba University}
\date{STOC 2027 submission \quad\NOTE{(author video for the PC)}}

\begin{document}
\maketitle

% Timing: 0:00--7:00
\section{The question}

% ============================================================
\begin{frame}{The calculus of relations (\CoR)}

\begin{center}
\begin{tikzpicture}[every node/.style={font=\small}]
  \node[cellvert,minimum size=6pt,label=below:$u$] (x) at (0,0) {};
  \node[cellvert,minimum size=6pt,label=above:$w$] (z) at (1.4,.9) {};
  \node[cellvert,minimum size=6pt,label=below:$v$] (y) at (2.8,0) {};
  \draw[-{Stealth},thick] (x) -- node[above left] {$X$} (z);
  \draw[-{Stealth},thick] (z) -- node[above right] {$Y$} (y);
  \draw[-{Stealth},line width=1.5pt,black!40] (3.5,.3) -- (4.7,.3);
  \node[cellvert,minimum size=6pt,label=below:$u$] (x2) at (5.4,0) {};
  \node[cellvert,minimum size=6pt,label=below:$v$] (y2) at (8.2,0) {};
  \draw[-{Stealth},thick] (x2) -- node[above] {$X \mathbin{;} Y$} (y2);
\end{tikzpicture}
\end{center}
\[
  X \mathbin{;} Y \;\defeq\; \set{(u, v) : (u, w) \in X \text{ and } (w, v) \in Y \text{ for some } w}
\]

\vfill
\pause
Terms over binary relation symbols, built from
\[
  \cup, \quad \cap, \quad \overline{\,\cdot\,} \ \text{\NOTE{(complement)}}, \quad {}^{\conv} \ \text{\NOTE{(converse)}}, \quad
  \zero, \ \one, \ \id, \quad \text{and} \quad \EMPH{$;$}
\]
\NOTE{Every intermediate result is a \emph{binary} relation.}

\vfill
\pause
\begin{columns}[T]
\begin{column}{.48\textwidth}
\begin{block}{As a matrix program}
  $;$ \ = Boolean matrix product\\
  ${}^{\conv}$ = transpose, \ others entrywise
\end{block}
\end{column}
\begin{column}{.48\textwidth}
\begin{block}{As a query plan}
  $;$ \ = join followed by projection\\
  \NOTE{\fcite{Codd1970}, \fcite{brijder2019tods}}
\end{block}
\end{column}
\end{columns}

\note{
- CoR is the variable-free algebra of binary relations.
- Composition X;Y keeps the pairs (u,v) connected through some middle vertex w;
  the middle vertex is eliminated.
- Read as Boolean matrices, composition is matrix multiplication; converse is transpose;
  the other operations are entrywise.
- So a CoR term is a program over whole matrices, and every intermediate result is a matrix.
}
\end{frame}

% ============================================================
\begin{frame}{Three-variable logic and \CoR}

\begin{columns}[c]
\begin{column}{.40\textwidth}
\centering
\begin{tikzpicture}[every node/.style={font=\small}]
  \node[cellvert,minimum size=7pt,label=left:$x$] (x) at (0,0) {};
  \node[cellvert,minimum size=7pt,label=right:$y$] (y) at (2.4,0) {};
  \node[cellvert,minimum size=7pt,label=above:$z$] (z) at (1.2,1.6) {};
  \draw[thick] (x) -- (y) -- (z) -- (x);
\end{tikzpicture}

\vspace{.3em}
\emph{\FOthree}: first-order logic with\\the variables $x, y, z$ only
\end{column}
\begin{column}{.16\textwidth}
\centering
\Large $\equiv$
\end{column}
\begin{column}{.40\textwidth}
\centering
\begin{tikzpicture}[every node/.style={font=\small}]
  \foreach \i/\x in {0/0,1/1.2,2/2.4}{\node[cellvert,minimum size=7pt] (a\i) at (\x,.8) {};}
  \draw[-{Stealth},thick] (a0) -- (a1);
  \draw[-{Stealth},thick] (a1) -- (a2);
  \draw[-{Stealth},thick] (a0) to[bend right=40] (a2);
\end{tikzpicture}

\vspace{.3em}
\emph{\CoR}: binary relations only
\end{column}
\end{columns}

\vfill
\begin{theorem}[Tarski; \ \fcite{tarski-givant1987}]
  \FOthree{} formulas $\varphi(x, y)$ and \CoR{} terms define the same binary queries.
\end{theorem}

\vfill
\pause
\begin{center}
\renewcommand{\arraystretch}{1.3}
\begin{tabular}{@{}lll@{}}
  $\CoR \to \FOthree$ & size $O(|T|)$ & \NOTE{composition $\mapsto$ $\exists z$} \\
  $\FOthree \to \CoR$ & size \EMPH{$2^{O(|\varphi|)}$} & \NOTE{classical translation}
\end{tabular}
\end{center}

\note{
- An FO3 formula can keep three vertices in scope and constrain all three edges among them.
- A CoR term sees only binary relations.
- Tarski showed that nevertheless they define exactly the same binary queries.
- The translation from CoR to FO3 is linear; the classical one back is single exponential.
}
\end{frame}

% ============================================================
\begin{frame}{Why translate: fast evaluation \ \NOTE{\fcite{nakamura-uezato2027}}}

\begin{center}
\begin{tikzpicture}[every node/.style={font=\small}]
  \node[logicrow] (k4) at (0,1.2) {$\mathrm{FO}^{k}$, \ $k \ge 4$};
  \node[logicrow] (k3) at (0,0) {\FOthree};
  \node[logicrow] (k2) at (0,-1.2) {$\mathrm{FO}^{2}$};
  \node[anchor=west] at (2,1.2) {no $O(m^{k-\varepsilon})$ algorithm \ \NOTE{(under Sparse MAX-3-SAT)}};
  \node[anchor=west] at (2,0) {$\xrightarrow{\ \text{translate}\ }$ \ \CoR{} \ $\xrightarrow{\ \text{matrix mult.}\ }$ \
    \EMPH{$2^{O(|\varphi|)}$} $\cdot\, n^{\omega}$};
  \node[anchor=west] at (2,-1.2) {$O(|\varphi| \cdot m)$};
\end{tikzpicture}
\end{center}
\NOTE{\hfill $n$: domain size, \ $m$: number of tuples, \ naive evaluation of $\mathrm{FO}^k$: \ $O(|\varphi| \cdot n^k)$}

\vfill
\pause
\begin{center}
  A \CoR{} term $T$ is evaluated in $O(|T| \cdot n^{\omega})$:\\
  one matrix product per composition.
\end{center}

\vfill
\pause
\begin{alertblock}{}
  \centering
  The factor $2^{O(|\varphi|)}$ comes only from the translation $\FOthree \to \CoR$.\\
  Is it necessary?
\end{alertblock}

\note{
- Why do we care about this translation? Because it gives fast evaluation.
- In our PODS 2027 paper: with four or more variables, the naive bound cannot be improved
  under a fine-grained hypothesis. With three variables, translating into CoR and multiplying
  matrices beats the naive n^3.
- The exponential factor is entirely due to the translation. So: is it necessary?
}
\end{frame}

% ============================================================
\begin{frame}{Where the exponential comes from}

\[
  G_k(x, y) \defeq \exists z \bigwedge_{i=1}^{k}
  \bigl[(P_i(x, y) \land Q_i(x, z)) \lor (Q_i(x, y) \land P_i(x, z))\bigr]
  \qquad \text{\NOTE{size $O(k)$}}
\]

\vfill
\begin{columns}[c]
\begin{column}{.42\textwidth}
\centering
\begin{tikzpicture}[every node/.style={font=\small}]
  \node[cellvert,minimum size=7pt,label=left:$x$] (x) at (0,0) {};
  \node[cellvert,minimum size=7pt,label=right:$y$] (y) at (3,0) {};
  \node[cellvert,minimum size=7pt,label=above:$z$] (z) at (1.2,1.8) {};
  \draw[thick] (x) -- node[below] {read edge: profile $A$} (y);
  \draw[thick] (x) -- node[above left] {$A^c$} (z);
  \node at (1.0,2.5) {witness edge};
\end{tikzpicture}

\NOTE{profile of an edge $=$ $\set{i : P_i \text{ holds}}$\\
(on inputs where exactly one of $P_i, Q_i$ holds)}
\end{column}
\pause
\begin{column}{.54\textwidth}
The witness edge must have the\\
\emph{complementary} profile of the read edge.

\vspace{.5em}
Composition accesses the witness edge separately,\\
so the classical translation \emph{fixes} the read profile:
\[
  G_k \;\equiv\; \bigcup_{A \subseteq [k]} \bigl[\, H_A \cap (H_{A^c} \mathbin{;} \one) \,\bigr]
\]
\NOTE{$H_A \defeq \bigcap_{i \in A} P_i \cap \bigcap_{i \notin A} Q_i$} \qquad
\EMPH{$2^k$ compositions}
\end{column}
\end{columns}

\note{
- Think of each edge as having a k-bit profile: the set of i with P_i.
- G_k asks for a witness z whose edge from x has exactly the complementary profile of the edge (x,y).
- In FO3 this is a size-O(k) formula, because x, y, z are all in scope and the clauses compare the two edges coordinatewise.
- In CoR, the witness edge is accessed only through a composition, separately from the read edge.
  The classical translation therefore enumerates the 2^k possible read profiles.
}
\end{frame}

% ============================================================
\begin{frame}{Is the blow-up necessary?}

\begin{center}
\renewcommand{\arraystretch}{1.5}
\begin{tabular}{@{}lll@{}}
  \toprule
  target & complement & exponential lower bound \\
  \midrule
  positive \CoR{} & none & known \ \fcite{nakamura2022} \\
  $\Sigma_1$\CoR{} & on inputs only & known \ \fcite{nakamura2022} \\
  \textbf{full \CoR{}} & anywhere & \only<1>{open}\only<2->{\textbf{this work}} \\
  \bottomrule
\end{tabular}
\end{center}

\vfill
\NOTE{Full \CoR{} = Boolean matrix programs with arbitrary Boolean operations, transpose, and products.}

\note{
- Previously, exponential lower bounds were known when complement does not appear above a composition.
- The question for full CoR, i.e., arbitrary Boolean matrix programs, was open.
}
\end{frame}

% Timing: 5:00--8:00   (paper: Sec. 1.1)
% One slide: the hard formula G_k, its DNF, and where the lower bound sits.
% (The DNF idea was already shown on the 'unavoidable?' slide; no section of its own.)

% ============================================================
\begin{frame}{The hard formula $G_k$}

\[
  G_k(x, y) \defeq \exists z \bigwedge_{i=1}^{k}
  \bigl[(P_i(x, y) \land Q_i(x, z)) \lor (Q_i(x, y) \land P_i(x, z))\bigr]
  \qquad \text{\NOTE{($|G_k| = O(k)$)}}
\]

\begin{center}
\begin{tikzpicture}[every node/.style={font=\small}]
  \node[pv,minimum size=7pt,label=left:$x$] (x) at (0,0) {};
  \node[pv,minimum size=7pt,label=right:$y$] (y) at (2.6,0) {};
  \node[pv,minimum size=7pt,label=above:$z$] (z) at (1.6,.6) {};
  \draw[thick,cA] (x) -- (y);
  \draw[thick,cB] (x) -- (z);
  \node[anchor=west,align=left] at (4.2,.6)
    {at every $i$: one of the edges $xy$, $xz$ is in $P_i$,\\
     the other is in $Q_i$};
\end{tikzpicture}
\end{center}

\vspace{1em}
DNF: \ $G_k \equiv \bigcup_{A \subseteq [k]} \bigl[ H_A \cap (H_{A^c} \mathbin{;} \one) \bigr]$
\ : \ \EMPH{$2^k$ compositions}
\ \NOTE{($H_A \defeq \bigcap_{i \in A} P_i \cap \bigcap_{i \notin A} Q_i$)}
\\ \hspace*{2.6em} $\Rightarrow$ optimal up to a factor of $2$: \ every equivalent \CoR{} term needs $\ge 2^{k-1}$ compositions


\note{
\textbf{スライド概要}：$G_k$ は DNF で合成 $2^k$ 個、これは 2 倍を除いて最適。

\textbf{メモ}：
\begin{itemize}
  \item 各 $i$ で、辺 $xy$ と辺 $xz$ の一方が $P_i$、他方が $Q_i$
  \item 下界は合成以外の演算（補・逆・共有）を無料にしても成り立つ
\end{itemize}
}
\end{frame}

% Timing: 10:00--13:00   (paper: Sec. 1.2)
\section{Natural approaches that fail}

% ============================================================
\begin{frame}{$\Sigma_1$\CoR{}: a fooling-set argument \ \NOTE{\fcite{nakamura2020}, \fcite{nakamura2022}}}

\begin{center}
\begin{tikzpicture}[every node/.style={font=\small}]
  % small structures
  \node[structbox,minimum width=2.6cm,minimum height=1.3cm] (Si) at (0,1.4) {};
  \node[above] at (Si.north) {$\mathcal S_i$};
  \node[pv,label=left:$u_i$] (ui) at (-.9,1.2) {};
  \node[pv,label=above:$w_i$] (wi) at (0,1.75) {};
  \node[pv,label=right:$v_i$] (vi) at (.9,1.2) {};
  \draw[thick,-{Stealth}] (ui) -- (wi); \draw[thick,-{Stealth}] (wi) -- (vi);
  \node[structbox,minimum width=2.6cm,minimum height=1.3cm] (Sj) at (0,-1.4) {};
  \node[below] at (Sj.south) {$\mathcal S_j$};
  \node[pv,label=left:$u_j$] (uj) at (-.9,-1.6) {};
  \node[pv,label=below:$w_j$] (wj) at (0,-1.05) {};
  \node[pv,label=right:$v_j$] (vj) at (.9,-1.6) {};
  \draw[thick,-{Stealth}] (uj) -- (wj); \draw[thick,-{Stealth}] (wj) -- (vj);
  \node[font=\scriptsize,text=chibasub] at (2.0,1.4) {$\varphi = 1$};
  \node[font=\scriptsize,text=chibasub] at (2.0,-1.4) {$\varphi = 1$};
  % large structure
  \begin{scope}[visible on=<2->]
  \node[structbox,minimum width=4.4cm,minimum height=2.4cm] (N) at (6.4,0) {};
  \node[above] at (N.north) {$\mathcal N$};
  \node[pv,label=left:$u^*$] (u) at (4.8,-.3) {};
  \node[pv,label=above:$c$] (c) at (6.4,.6) {};
  \node[pv,label=right:$v^*$] (v) at (8.0,-.3) {};
  \draw[thick,-{Stealth}] (u) -- (c); \draw[thick,-{Stealth}] (c) -- (v);
  \node[font=\scriptsize,text=chibasub] at (6.4,-.85) {$\varphi = 0$};
  \draw[-{Stealth},densely dashed,black!50] (Si.east) to[bend left=15] node[above,font=\scriptsize] {$\hookrightarrow$} (5.5,.4);
  \draw[-{Stealth},densely dashed,black!50] (Sj.east) to[bend right=15] node[below,font=\scriptsize] {$\hookrightarrow$} (7.3,.1);
  \end{scope}
\end{tikzpicture}
\end{center}

\vfill
\onslide<3->{%
\begin{center}
  $\Sigma_1$\CoR: \ $\iota(\sem{T}{\mathcal S}) \subseteq \sem{T}{\mathcal S'}$ \ $(\diamond)$
  \quad $\Rightarrow$ \quad each composition separates $\le 1$ instance
\end{center}}

\note{
\textbf{伝えること}：既存の指数下界は $\Sigma_1$\CoR{}（補が合成の上に来ない）だけ。その証明は fooling-set 論法として整理できる。
\begin{itemize}
  \item 小さい構造 $\mathcal S_i$ は $\varphi = 1$、大きい構造 $\mathcal N$ は $\varphi = 0$、交差した組は $\mathcal N$ に埋め込める
  \item 拡張で保存される $(\diamond)$ ので、各合成は 1 つのインスタンスしか分離できない $\to$ $t$ 個必要
\end{itemize}
}
\end{frame}

% ============================================================
\begin{frame}{With complement, both ingredients fail}

\begin{center}
\begin{tikzpicture}[every node/.style={font=\small}]
  \node[structbox,minimum width=2.6cm,minimum height=1.6cm] (S) at (0,0) {};
  \node[above] at (S.north) {$\mathcal S$};
  \node[pv,label=left:$u$] (u) at (-.5,.1) {};
  \node[font=\scriptsize] at (0,-.5) {$\overline{R \mathbin{;} \one}$ holds};
  \draw[-{Stealth},thick] (1.7,0) -- node[above] {$\iota$} (2.9,0);
  \node[structbox,minimum width=2.6cm,minimum height=1.6cm] (S2) at (4.6,0) {};
  \node[above] at (S2.north) {$\mathcal S'$};
  \node[pv,label=left:$\iota(u)$] (u2) at (4.2,.1) {};
  \node[pv] (r) at (5.3,.45) {};
  \draw[thick,-{Stealth},chibasub] (u2) -- node[above,font=\scriptsize] {$R$} (r);
  \node[font=\scriptsize,text=chibasub] at (4.6,-.5) {$\overline{R \mathbin{;} \one}$ fails};
\end{tikzpicture}

1. complement reverses $(\diamond)$
\end{center}

\vfill
\pause
\begin{center}
  2. \ $\overline{(E \cap (E \mathbin{;} \overline{\id}))} \mathbin{;} \one$ \ separates every $\mathcal S_i$ from $\mathcal N$
  \ \NOTE{(two compositions)}
\end{center}

\vfill
\pause
\HP{Unlike the inclusion $(\diamond)$, an \EMPH{equality} survives complement.}

\note{
\textbf{伝えること}：補を許すと、この論法の 2 つの材料が両方壊れる。教訓：\emph{包含}は補で反転するが、\emph{等式}は補を通る。
\begin{itemize}
  \item 1. 補は $(\diamond)$ を逆向きにする（図）
  \item 2. 合成 2 つの項で全インスタンスを分離できてしまう
  \item 他の手法（formula size、EF ゲーム、rank）も直接は使えない
  \item 最後の一行が、我々の不変条件を等式にする動機
\end{itemize}
}
\end{frame}

% Timing: 13:00--25:00   (paper: Sec. 1.3)
\section{Technical overview}

% ============================================================
\begin{frame}{Proof outline (at high-level)}

\KEY{Goal Lemma} \ Every \CoR{} term equivalent to $G_k$ needs $2^{\Omega(k)}$ compositions.

\vfill
We construct structures $\M_0$ and $\M_1, \dots, \M_{2^k}$ such that
\begin{itemize}
\uncover<2->{%
  \item $G_k$ distinguishes all of $\M_1, \dots, \M_{2^k}$ from $\M_0$:
    \begin{center}
    \renewcommand{\arraystretch}{1.15}
    \begin{tabular}{@{}w{l}{.6cm}|w{c}{.78cm}|*{7}{w{c}{.78cm}}|w{c}{.9cm}|w{c}{1.25cm}@{}}
      & $\M_0$ & $\M_{1}$ & $\M_{2}$ & \multicolumn{4}{w{c}{\dimexpr 3.12cm+6\tabcolsep\relax}}{$\cdots$} & $\M_{2^k}$ & \#diff. & \uncover<6->{\#comp.\,$\ge$} \\
      \hline
      $G_k$ & T & \alert{F} & \alert{F} & \multicolumn{4}{w{c}{\dimexpr 3.12cm+6\tabcolsep\relax}}{$\cdots$} & \alert{F} & \alert{$2^k$} & \uncover<6->{$2^k/2$} \\
    \end{tabular}
    \end{center}
}
\uncover<3->{%
   \item 
\KEY{Key Lemma} \ If a \CoR{} term $t$ distinguishes $d$ structures among $\M_1, \dots, \M_{2^k}$ from $\M_0$,
then $t$ needs $\ge d/2$ compositions.
\uncover<4->{%
\begin{center}
\renewcommand{\arraystretch}{1.15}
\begin{tabular}{@{}w{l}{.6cm}|w{c}{.78cm}|*{7}{w{c}{.78cm}}|w{c}{.9cm}|w{c}{1.25cm}@{}}
  & $\M_0$ & $\M_{1}$ & $\M_{2}$ & $\M_{3}$ & $\M_{4}$ & $\M_{5}$ & $\cdots$ & $\M_{2^k}$ & \#diff. & \#comp.\,$\ge$ \\
  \hline
  $t$ & T & \alert{F} & T & \alert{F} & T & T & $\cdots$ & T & \alert{$2$} & $1$ \\
  $t'$ & F & \alert{T} & \alert{T} & \alert{T} & \alert{T} & F & $\cdots$ & F & \alert{$4$} & $2$ \\
\end{tabular}
\end{center}
}}
\end{itemize}

\vfill
\uncover<6->{\HP{if $t \equiv G_k$, then $t$ needs $\ge 2^k/2$ compositions.}}

\vfill
\uncover<7->{\large \KEY{TASK} How can we construct such $\M_0$ and $\M_1, \dots, \M_{2^k}$?}

\note{
\textbf{スライド概要}：$G_k$ が $\M_0$ で成り立ち、$\M_1, \dots, \M_{2^k}$ で成り立たない構造を作る。$d$ 個を $\M_0$ から見分ける項は、合成を $d/2$ 個以上必要とする。

\textbf{メモ}：
\begin{itemize}
  \item 表の T/F は、各 $\M_i$ に対応する点対での値（論文の test pair）
  \item 以降：$\M_0$ の作り方、$\M_i$ ＝ hub での塗り替え $\M_\lambda$、合成が新たに見分けるのは高々 2 個である理由
\end{itemize}
}
\end{frame}

% ============================================================
% The structure M_0 (edge-colored complete graph, two middle vertices, the hub)
\begin{frame}{Structures as colored complete graphs}

\KEY{View} structures as \KEY{complete graphs}, colored with $2^k$ colors $\mu \subseteq [k]$:
$P_i \iff i \in \mu$.
\\ \NOTE{$Q_i$ is used as the complement of $P_i$ (i.e., $i \notin \mu$): exactly one of $P_i, Q_i$ holds on each edge.}

\vfill
\begin{center}
\begin{tikzpicture}[every node/.style={font=\small}]
  % same style as the figure of G_k (example with k = 3): faint candidates for z with dashed edges
  % (the dashed edges are drawn only near the candidates, to leave room for the label of x-z)
  \node[pv,minimum size=7pt,label=left:$x$] (x) at (0,0) {};
  \node[pv,minimum size=7pt,label=right:$y$] (y) at (4.8,0) {};
  \foreach \h in {2.45,3.0} {
    \node[pv,minimum size=5pt,draw=black!30] (c) at (3.0,\h) {};
    \draw[densely dashed,black!30,shorten <=2.2cm] (x) -- (c);
  }
  \node[font=\footnotesize,text=black!60] at (3.45,2.72) {$\exists z$};
  \node[pv,minimum size=7pt,label=right:$z$] (z) at (3.0,1.9) {};
  \draw[thick,cA] (x) -- node[below,text=cA]
    {color $\mu = \set{1, 3}$: \ $P_1, \ Q_2, \ P_3$} (y);
  \draw[thick,cB] (x) -- node[pos=.45,above left,align=right,text=cB,font=\footnotesize]
    {color $\mu^c = \set{2}$:\\ $Q_1, \ P_2, \ Q_3$} (z);
  \node[anchor=west,align=left] at (5.8,1.0)
    {for every $i \in [3]$,\\
    the $x$-$y$ and $x$-$z$ edges\\
    have labels $P_i$ and $Q_i$, in either order\\
    $\Rightarrow$ \ $G_3(x, y)$ \ \NOTE{($k = 3$)}};
\end{tikzpicture}
\end{center}

\vfill
\HP{$G_k(x, y)$ holds iff some $x$-$z$ edge has the color $\mu^c$, where $\mu$ is the color of the $x$-$y$ edge.}

\note{
\textbf{スライド概要}：構造を色付き完全グラフとして見る。$G_k$ は「補の色の $x$-$z$ 辺がある」と読める。

\textbf{メモ}：
\begin{itemize}
  \item 各辺で $P_i$ と $Q_i$ のちょうど一方が成り立つ
  \item 例：$k = 3$、$\set{1,3}$ と $\set{2}$ は互いに補
\end{itemize}
}
\end{frame}

% ============================================================
\begin{frame}{The randomly colored structure $\M_0$}

\begin{columns}[c]
  \begin{column}{.36\textwidth}
    \centering
    % Overall shape of M_0 (used on the slide 'The randomly colored structure M_0').
% Edge colors are fixed pseudo-random choices (seed 7); u (left), v (right) and the
% two middle vertices z (upper), z' (lower) of the pattern u -mu- z -nu- v are highlighted.
\begin{tikzpicture}[every node/.style={font=\scriptsize}]
  \foreach \i in {0,...,15} { \node[circle,fill=black,inner sep=0pt,minimum size=2.5pt] (m\i) at (180+360/16*\i:1.6) {}; }
  \draw[cC,line width=.25pt,opacity=.55] (m0) -- (m1);
  \draw[cB,line width=.25pt,opacity=.55] (m0) -- (m2);
  \draw[cD,line width=.25pt,opacity=.55] (m0) -- (m4);
  \draw[cA,line width=.25pt,opacity=.55] (m0) -- (m5);
  \draw[cA,line width=.25pt,opacity=.55] (m0) -- (m6);
  \draw[cA,line width=.25pt,opacity=.55] (m0) -- (m7);
  \draw[cC,line width=.25pt,opacity=.55] (m0) -- (m8);
  \draw[cA,line width=.25pt,opacity=.55] (m0) -- (m9);
  \draw[cB,line width=.25pt,opacity=.55] (m0) -- (m10);
  \draw[cA,line width=.25pt,opacity=.55] (m0) -- (m12);
  \draw[cA,line width=.25pt,opacity=.55] (m0) -- (m13);
  \draw[cD,line width=.25pt,opacity=.55] (m0) -- (m14);
  \draw[cD,line width=.25pt,opacity=.55] (m0) -- (m15);
  \draw[cA,line width=.25pt,opacity=.55] (m1) -- (m2);
  \draw[cB,line width=.25pt,opacity=.55] (m1) -- (m3);
  \draw[cA,line width=.25pt,opacity=.55] (m1) -- (m4);
  \draw[cD,line width=.25pt,opacity=.55] (m1) -- (m5);
  \draw[cA,line width=.25pt,opacity=.55] (m1) -- (m6);
  \draw[cA,line width=.25pt,opacity=.55] (m1) -- (m7);
  \draw[cB,line width=.25pt,opacity=.55] (m1) -- (m8);
  \draw[cA,line width=.25pt,opacity=.55] (m1) -- (m9);
  \draw[cD,line width=.25pt,opacity=.55] (m1) -- (m10);
  \draw[cA,line width=.25pt,opacity=.55] (m1) -- (m11);
  \draw[cB,line width=.25pt,opacity=.55] (m1) -- (m12);
  \draw[cA,line width=.25pt,opacity=.55] (m1) -- (m13);
  \draw[cB,line width=.25pt,opacity=.55] (m1) -- (m14);
  \draw[cC,line width=.25pt,opacity=.55] (m1) -- (m15);
  \draw[cD,line width=.25pt,opacity=.55] (m2) -- (m3);
  \draw[cB,line width=.25pt,opacity=.55] (m2) -- (m4);
  \draw[cA,line width=.25pt,opacity=.55] (m2) -- (m5);
  \draw[cC,line width=.25pt,opacity=.55] (m2) -- (m6);
  \draw[cB,line width=.25pt,opacity=.55] (m2) -- (m7);
  \draw[cA,line width=.25pt,opacity=.55] (m2) -- (m8);
  \draw[cB,line width=.25pt,opacity=.55] (m2) -- (m9);
  \draw[cC,line width=.25pt,opacity=.55] (m2) -- (m10);
  \draw[cA,line width=.25pt,opacity=.55] (m2) -- (m11);
  \draw[cA,line width=.25pt,opacity=.55] (m2) -- (m12);
  \draw[cA,line width=.25pt,opacity=.55] (m2) -- (m13);
  \draw[cB,line width=.25pt,opacity=.55] (m2) -- (m14);
  \draw[cD,line width=.25pt,opacity=.55] (m2) -- (m15);
  \draw[cD,line width=.25pt,opacity=.55] (m3) -- (m4);
  \draw[cC,line width=.25pt,opacity=.55] (m3) -- (m5);
  \draw[cD,line width=.25pt,opacity=.55] (m3) -- (m6);
  \draw[cD,line width=.25pt,opacity=.55] (m3) -- (m7);
  \draw[cC,line width=.25pt,opacity=.55] (m3) -- (m9);
  \draw[cC,line width=.25pt,opacity=.55] (m3) -- (m10);
  \draw[cB,line width=.25pt,opacity=.55] (m3) -- (m11);
  \draw[cB,line width=.25pt,opacity=.55] (m3) -- (m12);
  \draw[cB,line width=.25pt,opacity=.55] (m3) -- (m13);
  \draw[cA,line width=.25pt,opacity=.55] (m3) -- (m14);
  \draw[cC,line width=.25pt,opacity=.55] (m3) -- (m15);
  \draw[cD,line width=.25pt,opacity=.55] (m4) -- (m5);
  \draw[cC,line width=.25pt,opacity=.55] (m4) -- (m6);
  \draw[cD,line width=.25pt,opacity=.55] (m4) -- (m7);
  \draw[cC,line width=.25pt,opacity=.55] (m4) -- (m8);
  \draw[cA,line width=.25pt,opacity=.55] (m4) -- (m9);
  \draw[cA,line width=.25pt,opacity=.55] (m4) -- (m10);
  \draw[cD,line width=.25pt,opacity=.55] (m4) -- (m11);
  \draw[cB,line width=.25pt,opacity=.55] (m4) -- (m12);
  \draw[cC,line width=.25pt,opacity=.55] (m4) -- (m13);
  \draw[cB,line width=.25pt,opacity=.55] (m4) -- (m14);
  \draw[cD,line width=.25pt,opacity=.55] (m4) -- (m15);
  \draw[cD,line width=.25pt,opacity=.55] (m5) -- (m6);
  \draw[cA,line width=.25pt,opacity=.55] (m5) -- (m7);
  \draw[cA,line width=.25pt,opacity=.55] (m5) -- (m8);
  \draw[cC,line width=.25pt,opacity=.55] (m5) -- (m9);
  \draw[cC,line width=.25pt,opacity=.55] (m5) -- (m10);
  \draw[cC,line width=.25pt,opacity=.55] (m5) -- (m11);
  \draw[cD,line width=.25pt,opacity=.55] (m5) -- (m12);
  \draw[cD,line width=.25pt,opacity=.55] (m5) -- (m13);
  \draw[cA,line width=.25pt,opacity=.55] (m5) -- (m14);
  \draw[cA,line width=.25pt,opacity=.55] (m5) -- (m15);
  \draw[cC,line width=.25pt,opacity=.55] (m6) -- (m7);
  \draw[cD,line width=.25pt,opacity=.55] (m6) -- (m8);
  \draw[cA,line width=.25pt,opacity=.55] (m6) -- (m9);
  \draw[cA,line width=.25pt,opacity=.55] (m6) -- (m10);
  \draw[cC,line width=.25pt,opacity=.55] (m6) -- (m11);
  \draw[cD,line width=.25pt,opacity=.55] (m6) -- (m12);
  \draw[cC,line width=.25pt,opacity=.55] (m6) -- (m13);
  \draw[cD,line width=.25pt,opacity=.55] (m6) -- (m14);
  \draw[cC,line width=.25pt,opacity=.55] (m6) -- (m15);
  \draw[cA,line width=.25pt,opacity=.55] (m7) -- (m8);
  \draw[cD,line width=.25pt,opacity=.55] (m7) -- (m9);
  \draw[cC,line width=.25pt,opacity=.55] (m7) -- (m10);
  \draw[cB,line width=.25pt,opacity=.55] (m7) -- (m11);
  \draw[cA,line width=.25pt,opacity=.55] (m7) -- (m12);
  \draw[cD,line width=.25pt,opacity=.55] (m7) -- (m13);
  \draw[cA,line width=.25pt,opacity=.55] (m7) -- (m14);
  \draw[cB,line width=.25pt,opacity=.55] (m7) -- (m15);
  \draw[cC,line width=.25pt,opacity=.55] (m8) -- (m9);
  \draw[cB,line width=.25pt,opacity=.55] (m8) -- (m10);
  \draw[cB,line width=.25pt,opacity=.55] (m8) -- (m12);
  \draw[cD,line width=.25pt,opacity=.55] (m8) -- (m13);
  \draw[cD,line width=.25pt,opacity=.55] (m8) -- (m14);
  \draw[cD,line width=.25pt,opacity=.55] (m8) -- (m15);
  \draw[cA,line width=.25pt,opacity=.55] (m9) -- (m10);
  \draw[cB,line width=.25pt,opacity=.55] (m9) -- (m11);
  \draw[cD,line width=.25pt,opacity=.55] (m9) -- (m12);
  \draw[cD,line width=.25pt,opacity=.55] (m9) -- (m13);
  \draw[cC,line width=.25pt,opacity=.55] (m9) -- (m14);
  \draw[cB,line width=.25pt,opacity=.55] (m9) -- (m15);
  \draw[cD,line width=.25pt,opacity=.55] (m10) -- (m11);
  \draw[cC,line width=.25pt,opacity=.55] (m10) -- (m12);
  \draw[cD,line width=.25pt,opacity=.55] (m10) -- (m13);
  \draw[cC,line width=.25pt,opacity=.55] (m10) -- (m14);
  \draw[cD,line width=.25pt,opacity=.55] (m10) -- (m15);
  \draw[cB,line width=.25pt,opacity=.55] (m11) -- (m12);
  \draw[cB,line width=.25pt,opacity=.55] (m11) -- (m13);
  \draw[cA,line width=.25pt,opacity=.55] (m11) -- (m14);
  \draw[cB,line width=.25pt,opacity=.55] (m11) -- (m15);
  \draw[cB,line width=.25pt,opacity=.55] (m12) -- (m13);
  \draw[cB,line width=.25pt,opacity=.55] (m12) -- (m14);
  \draw[cB,line width=.25pt,opacity=.55] (m12) -- (m15);
  \draw[cA,line width=.25pt,opacity=.55] (m13) -- (m14);
  \draw[cD,line width=.25pt,opacity=.55] (m13) -- (m15);
  \draw[cB,line width=.25pt,opacity=.55] (m14) -- (m15);
  \draw[cA,line width=1.3pt] (m0) -- (m11);
  \draw[cB,line width=1.3pt] (m8) -- (m11);
  \draw[cA,line width=1.3pt] (m0) -- (m3);
  \draw[cB,line width=1.3pt] (m3) -- (m8);
  \node[left=0pt] at (m0) {$u$};
  \node[right=0pt] at (m8) {$v$};
  \node[above right=-1pt] at (m11) {$z$};
  \node[below left=-1pt] at (m3) {$z'$};
\end{tikzpicture}


    \NOTE{$\M_0$: randomly colored, $O(k\,4^k)$ vertices \fcite{jmt1995}}
  \end{column}
  \begin{column}{.60\textwidth}
    \KEY{Two-witness property} \ For all $u \ne v$ and all colors $\textcolor{cA}{\boldsymbol{\mu}}, \textcolor{cB}{\boldsymbol{\nu}}$,
    there are at least two \emph{witnesses} $z$ such that
    the $u$-$z$ and $z$-$v$ edges have colors $\textcolor{cA}{\boldsymbol{\mu}}$ and $\textcolor{cB}{\boldsymbol{\nu}}$, respectively.
  \end{column}
\end{columns}

\vfill
\pause
\vspace{.3em}
\KEY{Every \CoR{} term $t$ is color-invariant on $\M_0$}: $\M_0 \models t(u, v) \iff \M_0 \models t(u', v')$
whenever the $u$-$v$ and $u'$-$v'$ edges have the same color.

\vspace{1em}
Proof by induction on $t$: the case, $t = t_1 \mathbin{;} t_2$, follows from the (two-)witness property.

\note{
\textbf{スライド概要}：$\M_0$ ではどの色の道のパターンにも中間の頂点がある。だから項は「色 $\mapsto$ T/F」の表になる。

\textbf{メモ}：
\begin{itemize}
  \item 例：$P_1 \mathbin{;} Q_1$ はどの辺でも T（色 $\set{1}$ と $\emptyset$ の道が必ずある）、ループでは F
  \item 「2 つ以上」が効くのは $\M_\lambda$ 側（hub を避ける中間の頂点が残る）
\end{itemize}
}
\end{frame}

% ============================================================
\begin{frame}{$\M_\lambda$: recoloring at the hub $\hub$}

% Figure 1 of the paper, redrawn around the mechanism: the read edge at the hub
% accepts only witness edges of color lambda, and the fold recolors exactly those.
% Colors as on the slide of G_k: read edge cA, witness edges cB.
\begin{center}
\begin{tikzpicture}[every node/.style={font=\small},x=1.1cm,y=1.1cm,
    xmark/.style={text=red,font=\large\bfseries,inner sep=0pt}]
  \foreach \sx/\name/\lab/\vis in {0/0/structure $\M_0$/1-,7.3/1/recolored structure $\M_\lambda$/2-} {
    \begin{scope}[xshift=\sx cm,visible on=<\vis>]
      \node at (.6,2.35) {\lab};
      \node[pv,label=below:$\hub$] (h\name) at (0,0) {};
      \node[pv,label=below:$v$] (v\name) at (2.4,0) {};
      \node[pv] (z\name) at (-.5,1.6) {};
      \node[pv] (w\name) at (.6,1.75) {};
      \node[pv] (o\name) at (-1.6,.3) {};
      \draw[black!35] (h\name) -- node[below,font=\scriptsize,text=black!50] {$\mu \ne \lambda$} (o\name);
      \draw[line width=1.4pt,cA] (h\name) -- node[below=2pt,font=\scriptsize] {color $\lambda^c$} (v\name);
    \end{scope}
  }
  % M_0: witness edges of color lambda at the hub
  \draw[thick,cB] (h0) -- node[left,font=\scriptsize] {$\lambda$} (z0);
  \draw[thick,cB] (h0) -- node[right,font=\scriptsize] {$\lambda$} (w0);
  \node[font=\scriptsize,text=cB] at (-.2,2.0) {color $\lambda$};
  \node at (1.2,-1.05) {$G_k(\hub, v) = 1$};
  % M_lambda: the same edges, recolored to lambda^c, are no longer witnesses
  \begin{scope}[visible on=<2->]
    \draw[thick,black!35,densely dashed] (h1) -- node[left,font=\scriptsize,text=black!60] {$\lambda^c$} (z1);
    \draw[thick,black!35,densely dashed] (h1) -- node[right,font=\scriptsize,text=black!60] {$\lambda^c$} (w1);
    \node[xmark] at ($(h1)!.75!(z1)$) {$\times$};
    \node[xmark] at ($(h1)!.75!(w1)$) {$\times$};
    \node at (7.3/1.1+1.2,-1.05) {$G_k(\hub, v) = \textcolor{red}{0}$};
    \draw[-{Stealth},thick] (3.3,1.0) -- node[above] {recolor $\lambda$} (4.6,1.0);
  \end{scope}
\end{tikzpicture}
\end{center}

\vfill
\begin{center}
  $G_k(\hub, v)$ needs a $\hub$-$z$ edge of color $\lambda$
  \onslide<2->{\\ recoloring $\lambda$ changes \emph{exactly those edges at $\hub$}, \ and nothing else}
\end{center}

\vfill
\onslide<3->{%
\begin{center}
  $2^k$ recolored structures $\M_\lambda$, one for each color $\lambda$ \ \NOTE{($=$ the $\M_1, \dots, \M_{2^k}$ of the outline)}
  \\ away from $\hub$, terms cannot distinguish $\M_\lambda$ from $\M_0$ \NOTE{(of the $\ge 2$ middle vertices, one avoids $\hub$)}
\end{center}}

\note{
\textbf{スライド概要}：塗り替えで $G_k$ が 1 から 0 に変わる。

\textbf{メモ}：
\begin{itemize}
  \item $\hub$ に接する色 $\lambda$ の辺だけ塗り替える
  \item 色ごとに 1 つ、計 $2^k$ 個
\end{itemize}
}
\end{frame}

% ============================================================
\begin{frame}{Why a composition newly distinguishes at most two $\M_\lambda$}

% Figure 2 of the paper, redrawn with one layout for all panels:
% the pair is the bottom edge, witnesses sit above (and below) it.
% Red with a cross: an edge recolored by the fold; thick green: a witness that survives.
\begin{center}
\begin{tikzpicture}[every node/.style={font=\small},x=1.2cm,y=1.2cm,
    xmark/.style={text=red,font=\large\bfseries,inner sep=0pt},
    lost/.style={thick,red,densely dashed},
    kept/.style={line width=1.4pt,cC}]
  % (a) M_0
  \begin{scope}
    \node at (.9,1.75) {(a) $\M_0$, at $(\hub, v)$};
    \node[pv,label=below:$\hub$] (h) at (0,0) {};
    \node[pv,label=above:$z$] (z) at (.9,1.0) {};
    \node[pv,label=below:$v$] (v) at (1.8,0) {};
    \draw[thick] (h) -- node[above left] {$\lambda$} (z);
    \draw[thick] (z) -- node[above right] {$\nu$} (v);
    \draw[densely dotted] (h) -- (v);
    \node[align=center] at (.9,-.85) {middle vertex $z$:\\ $(X \mathbin{;} Y)(\hub, v) = 1$};
  \end{scope}
  % (b) M_lambda, at the hub: every witness uses an edge at the hub
  \begin{scope}[xshift=4.4cm,visible on=<2->]
    \node at (.9,1.75) {(b) $\M_\lambda$, at $(\hub, v)$};
    \node[pv,label=below:$\hub$] (h) at (0,0) {};
    \node[pv,label=above:$z$] (z) at (.9,1.0) {};
    \node[pv,label=below:$v$] (v) at (1.8,0) {};
    \draw[lost] (h) -- node[above left] {$\lambda^c$} (z);
    \node[xmark] at ($(h)!.5!(z)$) {$\times$};
    \draw[thick] (z) -- node[above right] {$\nu$} (v);
    \draw[densely dotted] (h) -- (v);
    \node[align=center] at (.9,-.85) {every middle vertex is lost:\\ \alert{distinguished}};
  \end{scope}
  % (c) M_lambda, away from the hub: the hub is only one of the witnesses
  \begin{scope}[xshift=8.8cm,visible on=<3->]
    \node at (.9,1.75) {(c) $\M_\lambda$, at $(u, v)$};
    \node[pv,label=left:$u$] (u) at (0,0) {};
    \node[pv,label=above:$\hub$] (h) at (.9,1.0) {};
    \node[pv,label=right:$v$] (v) at (1.8,0) {};
    \node[pv,label=below:$z'$] (z) at (.9,-.8) {};
    \draw[lost] (u) -- node[above left] {$\lambda^c$} (h);
    \node[xmark] at ($(u)!.5!(h)$) {$\times$};
    \draw[thick] (h) -- node[above right] {$\nu$} (v);
    \draw[kept] (u) -- node[below left] {$\lambda$} (z);
    \draw[kept] (z) -- node[below right] {$\nu$} (v);
    \draw[densely dotted] (u) -- (v);
    \node[align=center] at (.9,-1.45) {a second middle vertex $z'$\\ avoids $\hub$: \ not distinguished};
  \end{scope}
\end{tikzpicture}
\end{center}
\NOTE{\hfill a composition $X \mathbin{;} Y$, \ where $X$ holds on the edges of color $\lambda$ and $Y$ on those of color $\nu \ne \lambda$}

\vfill
\onslide<4->{%
\begin{center}
  only pairs at the hub can change, and only if $X$ or $Y$ singles out the color $\lambda$
  \\ $\Rightarrow$ \ each composition newly distinguishes at most \alert{two} $\M_\lambda$ from $\M_0$ \ \NOTE{(hence the Key Lemma)}
\end{center}}

\note{
\textbf{スライド概要}：$\hub$ を含む点対では見分けられうる、含まない点対では見分けられない。

\textbf{メモ}：
\begin{itemize}
  \item $\hub$ を含まない点対では、$\hub$ を通らない中間の頂点 $z'$ が残る
  \item 条件は 3 つ、どれも $\lambda$ を決める $\to$ 高々 2
\end{itemize}
}
\end{frame}

% Timing: 27:00--28:00   (paper: Sec. 7)
\section{Conclusion}

% ============================================================
\begin{frame}{Conclusion}

\begin{theorem}
There are \FOthree{} formulas $\varphi_1, \varphi_2, \dots$ with $|\varphi_k| = O(k)$
such that every \CoR{} term equivalent to $\varphi_k$ needs $\ge 2^{k-1}$ compositions.
\end{theorem}
Our proof also applies to \FOthree{} vs.\ \FOtwo{}, matrix query languages, and \CoR{} with RTC.

\medskip
\KEY{Key Idea} \ the randomly colored structure $\M_0$ with its $2^k$ flipped structures.\\
Each composition can newly distinguish $\le 2$ of them.

\pause
\vfill
\begin{problem}[Open Questions]
\begin{itemize}
  \item a superpolynomial gap over a \emph{fixed} finite signature (e.g., one binary relation)?
  \pause
  \item the domain size from which the blowup begins? \ \NOTE{(known: between $2^k/(8k)$ and $O(k 4^k)$)}
  \pause
  \item the exact number of compositions for $G_k$: between $2^{k-1}$ and $2^{k}$?
\end{itemize}
\end{problem}

\vfill
\emph{Thank you for your attention!}

\note{
\textbf{ページ概要}：
\begin{itemize}
  \item 結果（合成 $\ge 2^{k-1}$）と鍵となるアイデア（$\M_0$ の $2^k$ 通りの flip）のまとめ
  \item 未解決問題：固定シグネチャ、爆発が始まる領域サイズ、$G_k$ の正確な合成数
\end{itemize}
}
\end{frame}


\appendix
\setbeamertemplate{page number in head/foot}{}
\AtBeginSection[]{}
\begingroup
\AtNextBibliography{\small}
\begin{frame}[allowframebreaks,noframenumbering]{References}
  \printbibliography[heading=none]
\end{frame}
\endgroup

\end{document}
